% !TeX root = ../../Book1.tex
% 纲要标题：### 3.1 循环群
% 纲要位置：计划纲要.md，第 192 行。

\section{循环群}
\label{b1:sec:0301}

% 纲要内容（仅作写作计划，尚非教材正文）：
%
% * 有限及无限循环群的分类
% * 子群、商群与同态的显式计算
% * 循环群的自同构群
%

循环群中的每个元素都是一个固定元素的整数次幂，因此群内的等式可以转化为指数之间的整除关系。从整数加法群的子群出发，第一同构定理同时给出有限循环群和无限循环群的模型。

\ProseParagraphBreak
分类完成以后，还要能实际计算：哪些幂生成同一个子群，商群有多大，生成元可以映到哪里。下面始终区分“选择一个生成元”和“这个生成元是否由群本身唯一确定”。

\subsection{循环群的分类}
\label{b1:subsec:0301:01}

\begin{lemma}\label{b1:lem:integer-subgroups}
\(\mathbb Z\) 的每个加法子群都唯一写成 \(n\mathbb Z\)，其中 \(n\geq0\) 为整数。约定 \(0\mathbb Z=\{0\}\)。
\end{lemma}
\begin{proof}
零子群对应 \(n=0\)。非零子群含一个正整数，取其中最小者 \(n\)。任给子群中的 \(a\)，除法算法给出 \(a=qn+r\)，\(0\leq r<n\)。因
\(r=a-qn\) 仍在子群，最小性迫使 \(r=0\)。故子群包含于 \(n\mathbb Z\)，反向包含由 \(n\) 属于子群得出。正生成元是该子群最小的正元素，因而唯一。
\end{proof}

\begin{theorem}[循环群分类]\label{b1:thm:cyclic-classification}
每个无限循环群都同构于 \((\mathbb Z,+)\)；每个 \(n\) 阶有限循环群都同构于 \((\mathbb Z/n\mathbb Z,+)\)。因此两个循环群同构当且仅当它们的阶相同，包括无限阶情形。
\end{theorem}
\begin{proof}
设 \(G=\langle g\rangle\)。映射 \(\pi:\mathbb Z\rightarrow G\)，\(k\mapsto g^k\)，是满同态。核是 \(\mathbb Z\) 的子群，故\cref{b1:lem:integer-subgroups}给出唯一 \(n\geq0\)，使 \(\ker\pi=n\mathbb Z\)。对这个像为全体 \(G\) 的同态应用\cref{b1:thm:first-isomorphism}，得到同构 \(\mathbb Z/n\mathbb Z\rightarrow G\)，\([k]\mapsto g^k\)。若 \(n=0\)，商就是 \(\mathbb Z\)；若 \(n>0\)，\(0,\ldots,n-1\) 为商的全部不同代表，所以群阶为 \(n\)。同构保持元素个数，因而也得到最后的判别。
\end{proof}

记 \(n\) 阶循环群的一个标准模型为 \(C_n\coloneqq\mathbb Z/n\mathbb Z\)，其中 \(n\geq1\)，\(C_1\) 是平凡群。需要强调乘法表示时，也写 \(C_n=\langle g\rangle\)、\(g^n=1\)，同时声明 \(g\) 的阶恰为 \(n\)。仅写 \(g^n=1\) 还不够，因为阶可能是 \(n\) 的真因子。
\symidx{\(C_n\)：\(n\) 阶循环群}

\ProseParagraphBreak
上述证明给出的同构依赖生成元 \(g\) 的选择。若把生成元换成 \(g^a\)，由整数指数得到的同构一般也改变。这与上一章典范同构的讨论一致：有时群的同构类型唯一，但某个具体同构并不由群本身唯一确定。

\subsection{子群、商群与同态的计算}
\label{b1:subsec:0301:02}

\begin{proposition}\label{b1:prop:power-order}
设 \(G\) 为群，\(g\in G\)。若 \(g\) 的有限阶为正整数 \(n\)，则对任意 \(a\in\mathbb Z\)，有
\[
\operatorname{ord}(g^a)=\frac{n}{\gcd(n,a)}.
\]
这里 \(\gcd(n,a)\) 取正值，特别地，\(\gcd(n,0)=n\)，所以公式也包括 \(g^0=1\) 的情形。若 \(g\) 为无限阶且 \(a\in\mathbb Z\setminus\{0\}\)，则 \(g^a\) 也为无限阶。
\end{proposition}
\begin{proof}
有限情形中，\((g^a)^k=g^{ak}\)，由\eqref{b1:eq:element-order-divisibility}，它等于 \(1\) 当且仅当 \(n\mid ak\)。写
\(n=dn'\)、\(a=da'\)，其中 \(d=\gcd(n,a)\)、\(\gcd(n',a')=1\)。整除条件等价于 \(n'\mid a'k\)，再由互素性等价于 \(n'\mid k\)，所以最小正指数为 \(n'\)。无限情形中若正整数 \(k\) 满足 \(g^{ak}=1\)，则 \(ak\ne0\)，与无限阶矛盾。
\end{proof}

互素时的整除步骤可直接用 Bézout 等式验证：若 \(un'+va'=1\)，则
\(k=un'k+va'k\)，右边两项都被 \(n'\) 整除。这样的整数等式可由 Euclid 除法反复回代获得，本节只使用这一初等算术事实。

\begin{theorem}\label{b1:thm:cyclic-subgroups}
设 \(G\) 为循环群。则 \(G\) 的每个子群以及关于任意正规子群的商群均为循环群。若 \(G=\langle g\rangle\) 的有限阶为正整数 \(n\)，则对每个正因子 \(d\mid n\)，恰有一个 \(d\) 阶子群，即 \(\langle g^{n/d}\rangle\)，\(G\) 关于这个子群的商群阶为 \(n/d\)。
\end{theorem}
\begin{proof}
对于子群 \(H\leq G\)，取满同态 \(\pi:\mathbb Z\rightarrow G\)，\(k\mapsto g^k\)。原像 \(\pi^{-1}(H)\) 对整数之差封闭且含零，因而是加法子群。\cref{b1:lem:integer-subgroups}给出 \(\pi^{-1}(H)=m\mathbb Z\)；由 \(\pi\) 满射，\(\pi\bigl(\pi^{-1}(H)\bigr)=H\)，故 \(H=\langle g^m\rangle\)，所以循环。有限情形中原像包含 \(n\mathbb Z\)，于是 \(m\mid n\)，由\cref{b1:prop:power-order}，\(\langle g^m\rangle\) 的阶为 \(n/\gcd(n,m)=n/m\)，由阶唯一确定 \(m\)。群 \(G\) 交换，所以每个子群正规；商群的每个元素都是 \(gH\) 的幂，故商群循环。在有限情形中，其阶由\cref{b1:thm:lagrange}给出，即 \(\lvert G/H\rvert=\lvert G\rvert/\lvert H\rvert\)。
\end{proof}

例如 \(C_{18}=\langle g\rangle\) 中 \(g^6\) 的阶为 \(3\)，\(g^4\) 的阶为 \(9\)，且
\(\langle g^4\rangle=\langle g^2\rangle\)。商去 \(\langle g^6\rangle\) 得到六阶循环群。此时“子群唯一”是指给定阶的子群唯一，并非生成元唯一：同一个三阶子群由 \(g^6\) 或 \(g^{12}\) 生成。

\ProseParagraphBreak
同态同样可以显式计算。
\begin{proposition}[循环群之间的同态]\label{b1:prop:cyclic-homomorphisms}
设 \(m,n\) 为正整数，\(C_m=\langle x\rangle\)、\(C_n=\langle y\rangle\)。任一同态由 \(x\) 的像 \(y^a\) 唯一决定；指定这个像能够给出同态，当且仅当 \(n\mid ma\)。若 \(d=\gcd(m,n)\)，则模 \(n\) 的全部可取指数为
\[
a=0,n/d,2n/d,\ldots,(d-1)n/d,
\]
所以同态共有 \(d\) 个。
\end{proposition}
\begin{proof}
同态必须把 \(x^k\) 送到 \(y^{ak}\)，故至多有一个。因为 \(x^m=1\)，其存在要求 \(y^{am}=1\)，由\eqref{b1:eq:element-order-divisibility}等价于 \(n\mid ma\)。反过来，若此整除条件成立，则 \(x^k=x^\ell\) 给出 \(m\mid k-\ell\)，进而 \(n\mid a(k-\ell)\)，所以 \(y^{ak}=y^{a\ell}\)。公式 \(x^k\mapsto y^{ak}\) 因而良定义；指数相加立即给出同态性。写 \(m=dm'\)、\(n=dn'\)，则 \(n\mid ma\) 等价于 \(n'\mid m'a\)。因 \(m',n'\) 互素，Bézout 等式使它等价于 \(n'\mid a\)，给出所列的 \(d\) 个模 \(n\) 指数。
\end{proof}

比如 \(C_6\rightarrow C_9\) 只有三种，由生成元分别送到 \(1,y^3,y^6\) 给出；其中两种非平凡同态的像都为同一个三阶子群。

\subsection{循环群的自同构群}
\label{b1:subsec:0301:03}

\begin{definition}\label{b1:def:automorphism}
设 \(G\) 为群。将群同态 \(G\rightarrow G\) 称为 \(G\) 的\term[zitongtai]{自同态}[endomorphism]，将群同构 \(G\rightarrow G\) 称为 \(G\) 的\term[zitonggou]{自同构}[automorphism]。所有自同构在复合下构成\term[zitonggou-qun]{自同构群}[automorphism group]，记为 \(\operatorname{Aut}(G)\)。
\end{definition}
\symidx{\(\operatorname{Aut}(G)\)：自同构群}

复合封闭来自同态复合与双射复合；单位元是恒等映射；逆映射仍为同构。这说明自同构本身也有可研究的代数结构，它描述了群内部允许的重新命名方式。

\begin{theorem}[循环群的自同构]\label{b1:thm:cyclic-automorphisms}
设 \(n\geq2\) 为整数，\(C_n=\langle g\rangle\)。对每个满足 \(\gcd(a,n)=1\) 的整数 \(a\)，令
\[
\alpha_a:C_n\longrightarrow C_n,\qquad
\alpha_a(g^k)=g^{ak}\quad(k\in\mathbb Z).
\]
映射 \(\alpha_a\) 只依赖 \(a\) 的模 \(n\) 剩余类，且这些映射恰为 \(C_n\) 的全部自同构。复合公式为 \(\alpha_a\alpha_b=\alpha_{ab}\)，所以
\(\operatorname{Aut}(C_n)\) 同构于模 \(n\) 的可逆剩余类在乘法下组成的群，记为 \((\mathbb Z/n\mathbb Z)^\times\)。
\end{theorem}
\symidx{\((\mathbb Z/n\mathbb Z)^\times\)：可逆剩余类群}
\begin{proof}
在\cref{b1:prop:cyclic-homomorphisms}中取定义域与目标的阶都为 \(n\)，得到全部自同态 \(g^k\mapsto g^{ak}\)，其中任意模 \(n\) 指数 \(a\) 都满足存在条件 \(n\mid na\)。它满射当且仅当 \(g^a\) 仍生成
\(C_n\)，由\cref{b1:prop:power-order}，\(\operatorname{ord}(g^a)=n/\gcd(n,a)\)，故这当且仅当 \(\gcd(a,n)=1\)；有限集合上满射即双射。若互素，Bézout 等式给出 \(ab\equiv1\pmod n\)，所以 \(\alpha_b\) 是逆映射。反之，剩余类有乘法逆时 \(ab-1\) 被 \(n\) 整除，任何公因子都整除 \(1\)，故互素。逐点计算
\(\alpha_a\bigl(\alpha_b(g^k)\bigr)=g^{abk}\)，得到复合公式及所述群同构。
\end{proof}

定理的证明为了计算而选定了生成元 \(g\)，但所得自同构就是 \(C_n\) 上的幂映射 \(x\mapsto x^a\)，所以映射本身不依赖生成元的选择。事实上，若改取生成元 \(h=g^u\)，则
\(\alpha_a(h^k)=\alpha_a(g^{uk})=g^{auk}=h^{ak}\)。因此，可逆剩余类与幂自同构之间的对应不依赖生成元；与之不同，前面由生成元 \(g\) 给出的模型同构 \([k]\mapsto g^k\) 一般会随 \(g\) 改变。

\ProseParagraphBreak
对正整数 \(n\)，将模 \(n\) 的可逆剩余类的个数记为 \(\varphi(n)\)，并称 \(\varphi\) 为\term[euler-hanshu]{Euler 函数}[Euler totient function]。例如 \(n=8\) 时可取 \(1,3,5,7\)，它们的平方模 \(8\) 都为 \(1\)。于是 \(\operatorname{Aut}(C_8)\) 有四个元素，但不是循环群，因为没有四阶元素。一个循环群的自同构群不一定循环。

\ProseParagraphBreak
无限循环群只有两个自同构：在 \(\mathbb Z\) 的加法模型中，同态必为
\(k\mapsto ak\)，满射当且仅当 \(a=1\) 或 \(a=-1\)。因此
\(\operatorname{Aut}(\mathbb Z)\) 是二阶群。有限与无限情形的共同原则是“生成元必须送到生成元”，差别只在于整数指数何时仍能生成全部元素。
